\documentclass[11pt,a4paper]{article}
\usepackage[utf8]{inputenc}
\usepackage[english,russian]{babel}
\usepackage{amsmath}
\usepackage{amsfonts}
\usepackage{booktabs}
\usepackage{hyperref}
\usepackage{listings}
\usepackage{xcolor}
\usepackage{geometry}

\geometry{margin=1in}

\definecolor{codegreen}{rgb}{0,0.6,0}
\definecolor{codegray}{rgb}{0.5,0.5,0.5}
\definecolor{backcolour}{rgb}{0.95,0.95,0.92}

\lstset{
    backgroundcolor=\color{backcolour},
    commentstyle=\color{codegreen},
    keywordstyle=\color{magenta},
    numberstyle=\tiny\color{codegray},
    stringstyle=\color{purple},
    basicstyle=\ttfamily\small,
    breakatwhitespace=false,
    breaklines=true,
    captionpos=b,
    keepspaces=true,
    numbers=left,
    numbersep=5pt,
    showspaces=false,
    showstringspaces=false,
    showtabs=false,
    tabsize=2,
    frame=single
}

\title{Технический отчет: Оптимизация Атмосферы Hosek-Wilkie \\ \large Интеграция Sky LUT и HDR-Tonemapping для \texttt{Clouds.hs}}
\author{Группа разработки графических систем}
\date{\today}

\begin{document}

\maketitle

\section{Проблема производительности при прямом расчете}
Модель Hosek-Wilkie содержит тяжелые трансцендентные функции ($\exp$, $\text{pow}$, $\text{acos}$). Вычисление этой модели напрямую внутри основного цикла маршинга облаков (\texttt{Clouds.hs}) для каждого шага луча драматически снижает производительность (ALU-bound). При 64 шагах маршинга для экрана 1080p это потребует миллиардов вычислений экспонент за кадр.

Решение заключается в разделении расчетов: модель Hosek-Wilkie вычисляется один раз за кадр в текстуру-таблицу **Sky LUT**, а шейдер облаков выполняет быструю выборку (\texttt{Texture Fetch}) за константное время $\mathcal{O}(1)$.

---

\section{Математика и структура Sky LUT}
Sky LUT представляет собой двухмерную текстуру (рекомендуемый размер $200 \times 200$ пикселей) в формате высокой точности \texttt{VK\_FORMAT\_R16G16B16A16\_SFLOAT}. Карта координат текстуры ($U, V$) кодирует сферическую геометрию купола неба.

\subsection{Координатная параметризация (UV Mapping)}
Чтобы эффективно распределить пиксели и избежать искажений у горизонта, координаты текстуры $U$ и $V$ отображаются на углы неба следующим образом:

\begin{equation}
    U = \frac{1 + \cos\gamma}{2} \iff \cos\gamma = 2U - 1
\end{equation}
\begin{equation}
    V = \sqrt{\max(0.0, \cos\theta)} \iff \cos\theta = V^2
\end{equation}

Где:
\begin{itemize}
    \item $\cos\gamma$ --- скалярное произведение вектора взгляда и вектора солнца (\texttt{dot(viewDir, sunDir)}).
    \item $\cos\theta$ --- вертикальный компонент вектора взгляда (\texttt{viewDir.y}). Использование квадратного корня при кодировании $V$ дает более высокую плотность пикселей вблизи горизонта, где градиенты атмосферы наиболее резкие.
\end{itemize}

---

\section{Математика HDR-Тонирования (Tonemapping)}
Физические модели атмосферы оперируют реальными значениями энергетической яркости (\textit{Radiance}), которые в полдень могут превышать значения $100.0$ единиц, в то время как стандартные мониторы отображают диапазон $[0.0, 1.0]$. Прямое отсечение (\texttt{clamp}) уничтожит детали и сделает края облаков «выжженными».

\subsection{1. Экспоненциальная коррекция экспозиции}
Перед тонированием выполняется масштабирование физической яркости $L_{\text{raw}}$ с помощью параметра экспозиции ($E$), который динамически адаптируется в зависимости от времени суток:
\begin{equation}
    L_{\text{exposed}} = L_{\text{raw}} \cdot e^{E}
\end{equation}

\subsection{2. Оператор ACES Filmic Tonemapping}
Для кинематографического сжатия динамического диапазона и сохранения насыщенности в полутонах применяется аппроксимация кривой ACES (Academy Color Encoding System):

\begin{equation}
    \text{ACES}(x) = \frac{x \cdot (ax + b)}{x \cdot (cx + d) + e}
\end{equation}

Где эмпирические коэффициенты для стандартизированного перевода в sRGB равны:\\
$a = 2.51$, $b = 0.03$, $c = 2.43$, $d = 0.59$, $e = 0.14$.

---

\section{Реализация в шейдерном коде Vulkan}

\subsection{Генерация Sky LUT (Compute или Fragment Шейдер)}
Этот шейдер выполняется один раз в начале кадра для заполнения LUT-текстуры.

\begin{lstlisting}[language=C, caption={Запекание физической модели в Sky LUT}]
layout(local_size_x = 16, local_size_y = 16) in;
layout(set = 0, binding = 0, rgba16f) writeonly image2D skyLutImage;

void main() {
    ivec2 texEL = ivec2(gl_GlobalInvocationID.xy);
    vec2 size = vec2(imageSize(skyLutImage));
    if (texEL.x >= size.x || texEL.y >= size.y) return;

    vec2 uv = vec2(texEL) / (size - 1.0);

    // Декодирование UV в физические углы
    float cosGamma = uv.x * 2.0 - 1.0;
    float cosTheta = uv.y * uv.y; // Восстановление из нелинейного V
    float sinTheta = sqrt(1.0 - cosTheta * cosTheta);
    
    // Восстановление фиктивного вектора направления взгляда
    vec3 viewDir = vec3(sinTheta, cosTheta, 0.0); 
    // Направление на солнце передается из Uniform Buffer
    
    vec3 radiance = evaluateHosekWilkieInternal(viewDir, sunDir, params);
    imageStore(skyLutImage, texEL, vec4(radiance, 1.0));
}
\end{lstlisting}

\subsection{Выборка из LUT внутри маршинга облаков (\texttt{Clouds.hs})}
Теперь шейдер облаков осуществляет моментальный доступ к предрассчитанным данным рассеивания.

\begin{lstlisting}[language=Haskell, caption={Выборка освещения неба в Haskell FIR DSL}]
-- Быстрая выборка рассеивания атмосферы из LUT
lookupSkyAmbient :: Vec3 Float -> Vec3 Float -> Program Pipeline (Vec3 Float)
lookupSkyAmbient viewDir sunDir = do
    let cosGamma = dot viewDir sunDir
    let cosTheta = viewDir.&Y -- Направление вверх (Зенит)
    
    -- Перевод в координаты LUT по формулам нелинейного распределения
    let u = (cosGamma + 1.0) * 0.5
    let v = sqrt (max 0.0 cosTheta)
    
    skyRadiance <- sample @"sky_lut_texture" (Vec2 u v)
    return (skyRadiance.&XYZ)
\end{lstlisting}

\subsection{Финальный вывод кадра с ACES Tonemapping}
Применяется в самом конце графического пайплайна (Post-Processing) к совмещенному буферу облаков и ландшафта.

\begin{lstlisting}[language=C, caption={Реализация ACES Tonemapping}]
vec3 applyAcesTonemap(vec3 rawLinearColor, float exposure) {
    // Применение экспозиции
    vec3 x = rawLinearColor * exp(exposure);
    
    // Математическая аппроксимация кривой ACES
    vec3 a = vec3(2.51);
    vec3 b = vec3(0.03);
    vec3 c = vec3(2.43);
    vec3 d = vec3(0.59);
    vec3 e = vec3(0.14);
    
    vec3 mapped = (x * (a * x + b)) / (x * (c * x + d) + e);
    return clamp(mapped, 0.0, 1.0);
}
\end{lstlisting}

---

\section{Сравнительный анализ архитектуры пайплайнов}

\begin{table}[h!]
\centering
\small
\begin{tabular}{llll}
\toprule
\textbf{Метрика} & \textbf{Прямой расчет (Инлайн)} & \textbf{Оптимизированный Sky LUT} \\ \midrule
\textbf{Сложность вычислений} & $\mathcal{O}(\text{Pixels} \times \text{Steps} \times \text{ALU}_{\text{heavy}})$ & $\mathcal{O}(\text{Pixels} \times \text{Steps} \times \text{Fetch}_{\text{fast}})$ \\
\textbf{Затраты VRAM} & 0 КБ & $\approx 156$ КБ ($200 \times 200$ RGBA16F) \\
\textbf{Качество заката (HDR)} & Склонность к артефактам \texttt{NaN} & Стабильное распределение энергии \\
\textbf{Нагрузка на регистры GPU} & Критическая (Register Spilling) & Минимальная (Один текстурный дескриптор) \\ \bottomrule
\end{tabular}
\caption{Оценка эффективности перехода на архитектуру Sky LUT.}
\end{table}

\end{document}

